% ch21_monads_beck.tex — Volume II Chapter 9

\chapter{Monads: Algebraic Theory and Beck's Theorem}

\begin{overviewbox}
Chapter 8 introduced monads and showed how every adjunction $F \dashv G$
gives rise to a monad $T = GF$ on $\mathcal{C}$. But it left a gap:
different adjunctions can produce the \emph{same} monad. Which adjunction
is the ``right'' one? And conversely, given a monad $(T,\eta,\mu)$, can
we always recover an adjunction that produces it?

The answer is yes --- and in fact there are two canonical adjunctions, one
minimal (the \term{Kleisli category} $\mathcal{C}_T$) and one maximal (the
\term{Eilenberg--Moore category} $\mathcal{C}^T$). All adjunctions that
resolve $T$ factor through these two extremes via a \term{comparison
functor}.

The climax of this chapter is \term{Beck's monadicity theorem}: a precise
criterion for when a functor $F : \mathcal{C} \to \mathcal{D}$ is
\emph{monadic} --- meaning $\mathcal{D}$ is equivalent, over $\mathcal{C}$,
to the Eilenberg--Moore category of the induced monad. Beck's theorem
transforms an abstract question about equivalence of categories into a
concrete checklist involving coequalizers, and is one of the most useful
recognition theorems in all of category theory.

Along the way we meet \term{$F$-split coequalizers}, a remarkable class
of coequalizers that ``know how to split themselves'' when lifted along
$F$, and we see why they are the right notion of coequalizer for algebraic
categories.
\end{overviewbox}


% ═══════════════════════════════════════════════════════════════════════════
\section{Recap: Monads from Adjunctions}
\label{sec:monads-from-adj}
% ═══════════════════════════════════════════════════════════════════════════

\subsection{The Induced Monad}

Let $F : \mathcal{C} \rightleftharpoons \mathcal{D} : G$ be an adjunction
$F \dashv G$ with unit $\eta : \mathrm{Id}_\mathcal{C} \Rightarrow GF$ and
counit $\varepsilon : FG \Rightarrow \mathrm{Id}_\mathcal{D}$.

Set $T = GF : \mathcal{C} \to \mathcal{C}$. The monad structure is:
\begin{itemize}
  \item \textbf{Unit:} $\eta : \mathrm{Id}_\mathcal{C} \Rightarrow T = GF$.
        This is already the adjunction unit.
  \item \textbf{Multiplication:} $\mu = G\varepsilon F : GFGF \Rightarrow GF$.
        At component $c$: $\mu_c = G(\varepsilon_{Fc}) : GF GF c \to GF c$.
\end{itemize}
The monad axioms --- associativity $\mu \circ \mu T = \mu \circ T \mu$ and
the unit laws $\mu \circ \eta T = \mathrm{id}_T = \mu \circ T\eta$ ---
follow directly from the triangle identities for $F \dashv G$.

\begin{dialog}
So the monad carries only the ``shadow'' of the adjunction that lands back
on $\mathcal{C}$. The functor $F$ disappears; only the composite $GF$
survives. This information loss is precisely what Beck's theorem measures.
\end{dialog}

\subsection{The Eilenberg--Moore Category}

Given a monad $(T,\eta,\mu)$ on $\mathcal{C}$, the
\term{Eilenberg--Moore category} $\mathcal{C}^T$ has:
\begin{itemize}
  \item \textbf{Objects:} \term{$T$-algebras}: pairs $(a, h)$ where
        $a \in \mathcal{C}$ and $h : Ta \to a$ (the \emph{structure map}),
        satisfying:
        \begin{align*}
          h \circ \eta_a &= \mathrm{id}_a & &\text{(unit law)} \\
          h \circ \mu_a &= h \circ Th & &\text{(associativity law)}
        \end{align*}
  \item \textbf{Morphisms:} $(a,h) \to (b,k)$: morphisms $f : a \to b$
        in $\mathcal{C}$ such that $k \circ Tf = f \circ h$.
\end{itemize}
The forgetful functor $G^T : \mathcal{C}^T \to \mathcal{C}$ sends $(a,h)
\mapsto a$ and has a left adjoint $F^T : \mathcal{C} \to \mathcal{C}^T$
sending $a \mapsto (Ta, \mu_a)$. The induced monad of $F^T \dashv G^T$ is
exactly $(T,\eta,\mu)$.

\subsection{The Kleisli Category}

The \term{Kleisli category} $\mathcal{C}_T$ has:
\begin{itemize}
  \item \textbf{Objects:} the same objects as $\mathcal{C}$.
  \item \textbf{Morphisms:} $a \to b$ in $\mathcal{C}_T$ are morphisms
        $a \to Tb$ in $\mathcal{C}$ (``$T$-decorated'' or effectful maps).
  \item \textbf{Composition:} if $f : a \to Tb$ and $g : b \to Tc$ in
        $\mathcal{C}$, then $g \star f = \mu_c \circ Tg \circ f : a \to Tc$.
  \item \textbf{Identity on $a$:} the unit $\eta_a : a \to Ta$.
\end{itemize}
There is a free functor $F_T : \mathcal{C} \to \mathcal{C}_T$ (identity
on objects, $f \mapsto \eta_b \circ f$ on morphisms) and a forgetful
functor $G_T : \mathcal{C}_T \to \mathcal{C}$ (objects to $Ta$, morphisms
$f$ to $\mu_b \circ Tf$). These form an adjunction $F_T \dashv G_T$
inducing $(T,\eta,\mu)$.

\begin{howtosay}
\begin{itemize}
  \item A morphism $f : a \to Tb$ in $\mathcal{C}$ is called a
        \term{Kleisli morphism} or an \term{effectful computation} from $a$
        to $b$.
  \item Kleisli composition is the categorical form of ``bind'' ($>>=$ in
        Haskell).
  \item $\mathcal{C}_T$ is the minimal resolution; $\mathcal{C}^T$ is the
        maximal resolution.
\end{itemize}
\end{howtosay}


% ═══════════════════════════════════════════════════════════════════════════
\section{The Comparison Functor}
\label{sec:comparison-functor}
% ═══════════════════════════════════════════════════════════════════════════

\subsection{Every Resolution Factors Between the Two Extremes}

Suppose $F \dashv G$ resolves the monad $T = GF$. We define a functor
\[
  K : \mathcal{D} \to \mathcal{C}^T
\]
called the \term{comparison functor}, by:
\begin{itemize}
  \item On objects: $d \mapsto (Gd, G\varepsilon_d)$.
  \item On morphisms: $f \mapsto Gf$.
\end{itemize}

\begin{lemma}
The pair $(Gd, G\varepsilon_d)$ is indeed a $T$-algebra: the map
$G\varepsilon_d : GFGd \to Gd$ satisfies the $T$-algebra axioms.
\end{lemma}

\begin{proof}
Unit law: $G\varepsilon_d \circ \eta_{Gd} = \mathrm{id}_{Gd}$ by the
triangle identity $G\varepsilon \circ \eta G = \mathrm{id}_G$.

Associativity: $G\varepsilon_d \circ \mu_{Gd} = G\varepsilon_d \circ
GFG\varepsilon_d$ is the equation $G\varepsilon_d \circ G(\varepsilon_{FGd})
= G\varepsilon_d \circ GF(G\varepsilon_d)$, which is the naturality of
$\varepsilon$ applied to $\varepsilon_d$.
\end{proof}

Furthermore, $K$ commutes with the forgetful functors: $G^T \circ K = G$.
The comparison functor measures ``how much of the structure of $\mathcal{D}$
is captured by $T$-algebras.''

\subsection{When is $K$ an Equivalence?}

The functor $F \dashv G$ is called \term{monadic} if the comparison functor
$K : \mathcal{D} \to \mathcal{C}^T$ is an equivalence of categories. In this
case, $\mathcal{D}$ \emph{is} (up to equivalence) the category of algebras
for the monad $T$ on $\mathcal{C}$.

Examples:
\begin{itemize}
  \item The free/forgetful adjunction for groups is monadic: $\mathbf{Grp}
        \simeq \mathbf{Set}^T$ for the free-group monad.
  \item The free/forgetful adjunction for rings is monadic.
  \item The global-sections/constant-sheaf adjunction for sheaves is
        \emph{not} monadic in general.
\end{itemize}

Beck's theorem gives the exact criterion.


% ═══════════════════════════════════════════════════════════════════════════
\section{Beck's Monadicity Theorem}
\label{sec:beck}
% ═══════════════════════════════════════════════════════════════════════════

\subsection{The Statement}

\begin{theorem}[Beck's Monadicity Theorem]
Let $F : \mathcal{C} \to \mathcal{D}$ be a functor with right adjoint
$G : \mathcal{D} \to \mathcal{C}$, and let $T = GF$ be the induced monad.
The comparison functor $K : \mathcal{D} \to \mathcal{C}^T$ is an
equivalence of categories if and only if:
\begin{enumerate}
  \item[(B1)] $F$ \term{reflects isomorphisms}: if $Ff$ is an iso in
              $\mathcal{C}$, then $f$ is an iso in $\mathcal{D}$.
  \item[(B2)] $\mathcal{D}$ has coequalizers of all \term{$F$-split
              pairs}, and $F$ preserves these coequalizers.
\end{enumerate}
\end{theorem}

\subsection{Proof Sketch via the Prooflog}

We prove that under (B1) and (B2), the comparison functor $K$ is full,
faithful, and essentially surjective.

\begin{proof}
We trace the essential surjectivity argument, as it is the most
illuminating. Let $(a, h : Ta \to a)$ be any $T$-algebra. We need to find
$d \in \mathcal{D}$ with $Kd \cong (a,h)$.

\begin{prooflog}
\proofstep{Consider the pair $Fa \underset{\eta_{Fa}}{\overset{F h}{\rightrightarrows}} FTa = F GFa$}{$h : Ta \to a$ in $\mathcal{C}$; apply $F$ to get $Fh : FTa \to Fa$ and also have $\varepsilon_{Fa} : FGFa \to Fa$}
\proofstep{Actually the relevant pair is $FTa \overset{F h}{\underset{\varepsilon_{Fa}}{\rightrightarrows}} Fa$}{Both $Fh$ and $\varepsilon_{Fa}$ are morphisms $FTa = FGFa \to Fa$; the pair $(Fh, \varepsilon_{Fa})$ is $F$-split}
\proofstep{This pair is $F$-split via: $h : Ta \to a$, $\eta_a : a \to Ta$}{Check: $G(Fh) \circ \eta_a = Th \circ \eta_a$; $G(\varepsilon_{Fa}) \circ \eta_a = \mathrm{id}_a$ (triangle identity)}
\proofstep{By (B2), the coequalizer $d = \mathrm{coeq}(Fh, \varepsilon_{Fa})$ exists in $\mathcal{D}$}{Let $q : Fa \to d$ be the coequalizer morphism}
\proofstep{Apply $G$: since $F$ preserves this coequalizer and $G$ is right adjoint, $Gq$ is a split epi in $\mathcal{C}$}{The $F$-split condition propagates through the adjunction}
\proofstep{One shows $Kd = (Gd, G\varepsilon_d) \cong (a, h)$ as $T$-algebras}{The algebra map $Gd \to a$ is constructed from $q$ via the universal property of the coequalizer; the $T$-algebra laws follow from the $F$-split splitting data}
\proofstep{Faithfulness of $K$ follows from $F$ reflecting isomorphisms (B1): $K$-faithful + $F$-iso-reflecting forces $K$-faithful}{If $Kf = Kg$ then $GFf = GFg$; since the coequalizer structure forces uniqueness, $f = g$}
\proofstep{Fullness follows from the coequalizer presentation of free algebras}{Every $T$-algebra morphism lifts uniquely to a morphism in $\mathcal{D}$ via the universal property}
\end{prooflog}

Therefore $K$ is an equivalence. The converse direction checks that if $K$
is an equivalence, then conditions (B1) and (B2) must hold by abstract
properties of equivalences and the structure of $\mathcal{C}^T$.
\end{proof}

\begin{dialog}
Why are $F$-split pairs the right notion? Because split coequalizers are
\emph{absolute}: they are preserved by every functor. So requiring the
coequalizer to exist only after applying $F$ is the minimal assumption
needed to do the construction. If you demand more --- say, all reflexive
coequalizers --- you get a stronger (``crude'') form of monadicity that is
easier to apply but less sharp.
\end{dialog}


% ═══════════════════════════════════════════════════════════════════════════
\section{$F$-Split Coequalizers}
\label{sec:f-split}
% ═══════════════════════════════════════════════════════════════════════════

\subsection{Definition}

\begin{definition}[$F$-split pair and split coequalizer]
A \term{split coequalizer} is a diagram
\[
  a \;\underset{g}{\overset{f}{\rightrightarrows}}\; b \;\overset{q}{\to}\; c
\]
together with morphisms $s : c \to b$ and $t : b \to a$ satisfying:
\begin{align*}
  q \circ s &= \mathrm{id}_c \\
  f \circ t &= \mathrm{id}_b \\
  g \circ t &= s \circ q
\end{align*}
Such a diagram is a coequalizer of $(f,g)$, and it is preserved by every
functor (the splitting data witnesses the coequalizer purely equationally).

A pair $(f, g) : a \rightrightarrows b$ in $\mathcal{D}$ is called an
\term{$F$-split pair} (for a functor $F : \mathcal{D} \to \mathcal{C}$)
if the pair $(Ff, Fg)$ has a split coequalizer in $\mathcal{C}$.
\end{definition}

\subsection{The Key Example: Free Algebra Pairs}

Given a $T$-algebra $(a, h : Ta \to a)$, the pair
\[
  T^2 a \;\underset{Th}{\overset{\mu_a}{\rightrightarrows}}\; Ta
\]
is always split in $\mathcal{C}$ (via $\eta_{Ta}$ and $\eta_a$), and its
coequalizer is $a$ with coequalizer map $h$. This is the prototype of
every $F$-split pair that arises in Beck's theorem.

\begin{example}{$F$-split pairs for the free-module monad}
Let $R$ be a ring, $F = R \otimes_\mathbb{Z} - : \mathbf{Ab} \to
R\text{-}\mathbf{Mod}$ the free $R$-module functor (left adjoint to
the forgetful functor $G$). A pair $(f,g) : M \rightrightarrows N$ in
$R\text{-}\mathbf{Mod}$ is $F$-split iff the underlying abelian group pair
$(Gf, Gg) : GM \rightrightarrows GN$ has a split coequalizer in
$\mathbf{Ab}$.

In practice: this pair is split iff the coequalizer in $\mathbf{Ab}$
happens to be the underlying abelian group of an $R$-module that $R$
can act on in a compatible way.
\end{example}

\subsection{Absolute Coequalizers}

The splitting data $(s,t)$ is remarkable: it makes the coequalizer
\emph{absolute}, meaning it is preserved by every functor without any
assumption on that functor. Contrast:
\begin{itemize}
  \item Ordinary coequalizers: preserved by left adjoints.
  \item Reflexive coequalizers: preserved by categories satisfying
        additional exactness properties.
  \item \textbf{Split coequalizers: preserved by every functor.}
\end{itemize}
This is why split coequalizers are the right setting for Beck's theorem:
the condition ``$F$ preserves the coequalizer'' is automatically satisfied
once the coequalizer is split.


% ═══════════════════════════════════════════════════════════════════════════
\section{Applications}
\label{sec:monads-applications}
% ═══════════════════════════════════════════════════════════════════════════

\subsection{Rings and Modules}

Let $R$ be a commutative ring. The free $R$-module functor
$F = R[-] : \mathbf{Set} \to R\text{-}\mathbf{Mod}$, sending a set $X$
to the free $R$-module $R[X] = \bigoplus_{x \in X} R$, is left adjoint to
the forgetful functor $G$.

\begin{theorem}
The adjunction $F \dashv G$ is monadic: $R\text{-}\mathbf{Mod} \simeq
\mathbf{Set}^T$ where $T = GF$ is the free $R$-module monad.
\end{theorem}

\begin{proof}
We verify Beck's conditions.
\begin{prooflog}
\proofstep{$G$ reflects isomorphisms}{A map of $R$-modules is an iso iff its underlying set map is a bijection; $G$ is the forgetful functor}
\proofstep{$G$ creates coequalizers of $G$-split pairs}{Given a $G$-split pair $f,g : M \rightrightarrows N$ in $R\text{-}\mathbf{Mod}$, its set-level split coequalizer has a unique $R$-module structure making it the coequalizer in $R\text{-}\mathbf{Mod}$}
\proofstep{$G$ preserves these coequalizers}{Since the coequalizer is split at the set level, it is preserved by every functor, including $G$}
\proofstep{Both Beck conditions satisfied; apply the theorem}{$K : R\text{-}\mathbf{Mod} \to \mathbf{Set}^T$ is an equivalence}
\end{prooflog}
\end{proof}

This result generalizes: for any algebraic theory (a set of operations and
equations), the category of models is equivalent to $\mathbf{Set}^T$ for
the corresponding monad.

\subsection{Algebraic Theories as Monads}

A \term{Lawvere theory} or \term{algebraic theory} is a category $\mathbf{T}$
with finite products, whose objects are finite powers $n = 1^n$ of a
distinguished object $1$. A model of $\mathbf{T}$ in $\mathbf{Set}$ is a
product-preserving functor $M : \mathbf{T} \to \mathbf{Set}$.

\textbf{Main fact}: the category of models $\mathrm{Mod}(\mathbf{T},
\mathbf{Set})$ is equivalent to the Eilenberg--Moore category $\mathbf{Set}^T$
for a naturally constructed monad $T$ on $\mathbf{Set}$. The monad $T(X)$
is the set of $\mathbf{T}$-terms in variables $X$.

This connection --- Beck's theorem applied to Lawvere theories --- is one
of the deepest results in categorical algebra: \emph{monads on $\mathbf{Set}$
are exactly algebraic theories} (in a precise sense, with appropriate
finiteness conditions).

\subsection{Crude Monadicity}

In practice, one often uses a simpler sufficient criterion:

\begin{theorem}[Crude Monadicity / Barr--Beck]
$F \dashv G$ is monadic if $F$ reflects isomorphisms and $\mathcal{D}$ has
and $G$ preserves coequalizers of all \emph{reflexive pairs} (pairs $(f,g)$
that admit a common right inverse).
\end{theorem}

This handles many algebraic cases without having to identify $F$-split pairs
explicitly.


% ═══════════════════════════════════════════════════════════════════════════
\section{Exercises}
\label{sec:monads-beck-exercises}
% ═══════════════════════════════════════════════════════════════════════════

\begin{exercisesbox}
\begin{qa}
\qaitem{What is the multiplication $\mu$ of the monad $T = GF$ arising from an adjunction $F \dashv G$?}{$\mu = G\varepsilon F : GFGF \Rightarrow GF$, where $\varepsilon$ is the counit of the adjunction. At component $c$: $\mu_c = G(\varepsilon_{Fc})$.}
\qaitem{What are the two canonical adjunctions that resolve a monad $T$?}{The Kleisli adjunction $F_T \dashv G_T$ with $\mathcal{C}_T$ (minimal, initial among all resolutions), and the Eilenberg--Moore adjunction $F^T \dashv G^T$ with $\mathcal{C}^T$ (maximal, terminal among resolutions).}
\qaitem{What is a $T$-algebra for a monad $(T,\eta,\mu)$?}{A pair $(a, h: Ta \to a)$ where $h \circ \eta_a = \mathrm{id}_a$ (unit law) and $h \circ \mu_a = h \circ Th$ (associativity law). The map $h$ is the ``structure map'' or ``action.'' }
\qaitem{What is the comparison functor $K$?}{Given an adjunction $F \dashv G$ inducing $T = GF$, the functor $K : \mathcal{D} \to \mathcal{C}^T$ sending $d \mapsto (Gd, G\varepsilon_d)$. It commutes with the forgetful functors over $\mathcal{C}$.}
\qaitem{State Beck's monadicity theorem.}{$K$ is an equivalence iff: (B1) $F$ reflects isomorphisms, and (B2) $\mathcal{D}$ has coequalizers of $F$-split pairs and $F$ preserves them.}
\qaitem{What is a split coequalizer?}{A diagram $a \rightrightarrows b \to c$ with extra maps $s: c \to b$ and $t: b \to a$ satisfying $qs = \mathrm{id}$, $ft = \mathrm{id}$, $gt = sq$. These equations make $c$ the coequalizer of $(f,g)$ without requiring any universal property argument.}
\qaitem{Why are split coequalizers preserved by every functor?}{Because the splitting is expressed purely by equations between morphisms, and every functor preserves equations. No colimit universal property needs to be checked.}
\qaitem{What does ``$F$-split pair'' mean?}{A pair $(f,g): a \rightrightarrows b$ in $\mathcal{D}$ such that $(Ff, Fg)$ has a split coequalizer in $\mathcal{C}$. The split structure may not lift to $\mathcal{D}$; it only needs to exist downstairs.}
\qaitem{Why is the free-module adjunction monadic?}{The forgetful functor $G: R\text{-}\mathbf{Mod} \to \mathbf{Set}$ reflects isomorphisms (module isos are exactly underlying-set bijections) and creates coequalizers of $G$-split pairs (the coequalizer set has a unique $R$-module structure).}
\qaitem{What is a morphism in the Kleisli category $\mathcal{C}_T$?}{A morphism $a \to b$ in $\mathcal{C}_T$ is a morphism $a \to Tb$ in $\mathcal{C}$. Kleisli composition of $f: a \to Tb$ and $g: b \to Tc$ is $\mu_c \circ Tg \circ f : a \to Tc$.}
\qaitem{What is the identity morphism on $a$ in $\mathcal{C}_T$?}{The unit $\eta_a : a \to Ta$.}
\qaitem{What does it mean for a functor to reflect isomorphisms?}{$F$ reflects isomorphisms if: whenever $Ff$ is an isomorphism in the codomain, $f$ was already an isomorphism in the domain.}
\qaitem{Give an example of an adjunction that is \emph{not} monadic.}{The inclusion $i: \mathbf{Ab} \hookrightarrow \mathbf{Grp}$ with left adjoint the abelianization functor $(-)/[G,G]$ is not monadic: the comparison functor is not full.}
\qaitem{What is the Crude Monadicity Theorem?}{$F \dashv G$ is monadic if $F$ reflects isomorphisms and the domain has and $G$ preserves coequalizers of reflexive pairs.}
\qaitem{What is a reflexive pair?}{A pair of morphisms $f, g: a \rightrightarrows b$ that share a common right inverse: there exists $s: b \to a$ with $fs = \mathrm{id}_b = gs$.}
\qaitem{Why does the Eilenberg--Moore category earn the name ``maximal resolution''?}{Every adjunction $F' \dashv G'$ resolving $T$ factors through $\mathcal{C}^T$ via the comparison functor $K$: $\mathcal{D} \xrightarrow{K} \mathcal{C}^T$. So $\mathcal{C}^T$ is final among all such factorizations.}
\qaitem{Why does the Kleisli category earn the name ``minimal resolution''?}{Every adjunction resolving $T$ receives a functor from $\mathcal{C}_T$: there is a canonical $J: \mathcal{C}_T \to \mathcal{D}$ for any resolution, making $\mathcal{C}_T$ initial.}
\qaitem{What is the free $T$-algebra on $a$?}{The pair $(Ta, \mu_a)$, where the structure map is $\mu_a : T^2 a \to Ta$. The associativity and unit laws for $(Ta, \mu_a)$ are exactly the monad axioms for $(T,\eta,\mu)$.}
\qaitem{What pair of morphisms does every $T$-algebra $(a,h)$ present as a coequalizer?}{The pair $\mu_a, Th : T^2 a \rightrightarrows Ta$. The coequalizer is $a$ with map $h: Ta \to a$. This pair is always split by $\eta_{Ta}$ and $\eta_a$.}
\qaitem{In functional programming, what corresponds to the Kleisli composition?}{The ``bind'' operation ($>>=$ in Haskell). If $f: a \to M\,b$ and $g: b \to M\,c$, then $g \star f = $ `a >>= f >>= g'. The monad laws for $>>=$ are exactly the Kleisli category axioms.}
\end{qa}
\end{exercisesbox}


% ═══════════════════════════════════════════════════════════════════════════
\begin{summarybox}
\textbf{The two canonical resolutions.}\quad
Every monad $(T,\eta,\mu)$ on $\mathcal{C}$ arises from an adjunction.
The minimal such adjunction lives in the Kleisli category $\mathcal{C}_T$
(objects of $\mathcal{C}$, morphisms are $T$-decorated maps); the maximal
one lives in the Eilenberg--Moore category $\mathcal{C}^T$ ($T$-algebras).
Every other resolution factors between these two extremes.

\textbf{The comparison functor.}\quad
Any adjunction $F \dashv G$ inducing $T$ gives a comparison functor
$K: \mathcal{D} \to \mathcal{C}^T$, $d \mapsto (Gd, G\varepsilon_d)$.
The question ``is this adjunction monadic?'' is the question ``is $K$ an
equivalence?''

\textbf{Beck's theorem.}\quad
$K$ is an equivalence iff $F$ reflects isomorphisms and $\mathcal{D}$
has, with $F$ preserving, coequalizers of $F$-split pairs. The key notion
is an $F$-split pair: one whose image under $F$ has a \emph{split}
(absolute) coequalizer.

\textbf{Algebraic significance.}\quad
The Eilenberg--Moore category of the free-algebra monad over $\mathbf{Set}$
is exactly the category of algebras for the corresponding algebraic theory.
This makes Beck's theorem the precise categorical statement that monads on
$\mathbf{Set}$ are the same as equational algebraic theories.
\end{summarybox}


% ═══════════════════════════════════════════════════════════════════════════
\section{Formal Restatement}
\label{sec:monads-beck-formal}
% ═══════════════════════════════════════════════════════════════════════════

\begin{formalrestatement}

\textbf{Monad.}\quad A monad on $\mathcal{C}$ is a monoid object in the
strict monoidal category $(\mathrm{End}(\mathcal{C}), \circ, \mathrm{Id})$
of endofunctors: a triple $(T, \eta: \mathrm{Id} \Rightarrow T,
\mu: T^2 \Rightarrow T)$ with $\mu \circ \mu T = \mu \circ T\mu$,
$\mu \circ \eta T = \mathrm{id}_T = \mu \circ T\eta$.

\textbf{Eilenberg--Moore category.}\quad $\mathcal{C}^T$ is the category
of $T$-algebras $(a, h: Ta \to a)$ satisfying the two coherence laws, with
morphisms the $T$-equivariant maps. The forgetful $G^T: \mathcal{C}^T \to
\mathcal{C}$ and free $F^T: \mathcal{C} \to \mathcal{C}^T$ ($a \mapsto
(Ta, \mu_a)$) form an adjunction inducing $T$.

\textbf{Kleisli category.}\quad $\mathcal{C}_T$ has the same objects as
$\mathcal{C}$ and $\mathrm{Hom}_{\mathcal{C}_T}(a,b) = \mathrm{Hom}_\mathcal{C}(a, Tb)$,
with Kleisli composition $g \star f = \mu \circ Tg \circ f$ and
identity $\eta_a$. The inclusion $F_T: \mathcal{C} \to \mathcal{C}_T$ and
resolution $G_T: \mathcal{C}_T \to \mathcal{C}$ form an adjunction inducing $T$.

\textbf{Comparison functor.}\quad For any adjunction $F \dashv G$ with
induced monad $T$, the comparison functor $K: \mathcal{D} \to \mathcal{C}^T$
is $Kd = (Gd, G\varepsilon_d)$, $Kf = Gf$. It satisfies $G^T \circ K = G$.

\textbf{Split coequalizer.}\quad A diagram $a \overset{f}{\underset{g}{\rightrightarrows}} b \overset{q}{\to} c$
with $s: c \to b$ and $t: b \to a$ satisfying $qs = \mathrm{id}_c$, $ft = \mathrm{id}_b$,
$gt = sq$. Every functor preserves split coequalizers.

\textbf{$F$-split pair.}\quad A pair $(f,g): a \rightrightarrows b$ in
$\mathcal{D}$ such that $(Ff, Fg)$ admits a split coequalizer in $\mathcal{C}$.

\textbf{Beck's Monadicity Theorem.}\quad $K$ is an equivalence of
categories iff (B1) $F$ reflects isomorphisms; and (B2) $\mathcal{D}$ has
coequalizers of all $F$-split pairs, and $F$ preserves them.

\textbf{Crude monadicity.}\quad Sufficient: (B1) holds and $\mathcal{D}$
has, with $G$ preserving, coequalizers of all reflexive pairs in $\mathcal{D}$.

\end{formalrestatement}
